\section{Introduction}
\label{sec:intro}

Vision-language models (VLMs), exemplified by CLIP~\cite{radford2021learning}, acquire broad visual concepts and open-vocabulary recognition capabilities by aligning images and natural language at scale. These properties make CLIP an appealing foundation for continual learning, where a shared model must absorb sequentially arriving classes without being retrained from scratch. Yet continual adaptation also changes what must be preserved. Beyond retaining discrimination over previously observed classes, an adapted CLIP model should maintain the image--text alignment that supports zero-shot recognition and cross-modal retrieval. We therefore study CLIP class-incremental learning through two complementary objectives: learning strong classifiers for sequentially introduced classes and preserving the cross-modal capability inherited from pre-training.

These objectives expose two related but distinct forms of degradation. Supervision from a new task reshapes the representation toward the currently observed classes and can interfere with directions that are important to historical tasks. Repeated adaptation can also move the representation away from the image--text relations learned during pre-training~\cite{zheng2023zscl,ni2023modx}. A frozen CLIP model retains its pre-trained geometry but cannot fully exploit new supervision, whereas unconstrained adaptation improves current-task discrimination at the cost of historical knowledge and cross-modal transfer. A continual learner must therefore control interference with historical task directions while explicitly anchoring cross-modal relations during training.

Low-rank adaptation (LoRA)~\cite{hu2022lora} provides a natural residual pathway for updating a frozen foundation model. For a linear transformation, LoRA parameterizes the residual update as $\Delta\mathbf{W}=\mathbf{A}\mathbf{B}$, with a rank much smaller than the ambient feature dimension. This factorization reduces the number of trainable parameters, but low rank alone is not a knowledge-preservation constraint: the input-side factor remains free to amplify directions on which historical tasks exhibit high activation energy. In other words, LoRA determines how many degrees of freedom are used for adaptation, but not which historical input directions these degrees of freedom should follow.

Null-space methods provide a complementary strategy by identifying directions that minimally affect previously acquired knowledge. Gradient-projection approaches modify an update after conventional forward and backward computation so that it lies in an approximate null space of historical features~\cite{saha2021gradient,kang2025dynamic}. Other null-space-inspired LoRA variants initialize or constrain the input-side factor using a basis derived from low-energy activation directions~\cite{tang2026lora}. These designs inject historical geometry into the gradient or a low-rank factor, but they do not place a fixed history-induced filter persistently before both trainable factors. This distinction motivates our central question: can historical activation geometry be embedded directly into the LoRA forward function, such that both low-rank factors are optimized in a filtered input space throughout adaptation?

We answer this question with \textbf{LoRA-NF}, low-rank adaptation with null-space filtering. Under the right-multiplication convention used throughout this paper, LoRA-NF replaces the conventional update $\mathbf{A}\mathbf{B}$ with
\begin{equation}
    \Delta\mathbf{W}=\mathbf{P}\mathbf{A}\mathbf{B},
    \label{eq:intro_lora_nf}
\end{equation}
where $\mathbf{P}$ is constructed from the eigenspectrum of historical layer-input second moments. The adapter consequently computes $(\mathbf{X}\mathbf{P})\mathbf{A}\mathbf{B}$: historical high-energy directions are attenuated before the two trainable factors learn their transformation. Unlike an initialization-only bias, the filter remains active throughout optimization. We use a leaky rather than strict filter, retaining an adaptation path in every direction while biasing learning away from directions most strongly occupied by historical tasks. The formal conditions governing its historical response, expressivity, and optimization dynamics are developed in Section~\ref{sec:lora_nf_theory}.

LoRA-NF is designed to preserve historical task knowledge; it is not, by itself, a mechanism for preserving cross-modal retrieval. We assign the latter role primarily to cross-modal knowledge distillation, which constrains the adapted model to reproduce the image--text relations of a frozen CLIP teacher on reference pairs. At inference, we further combine the CLIP text classifier with low-rank regularized Gaussian discriminant analysis (LR-RGDA), which captures class-conditional statistics in the adapted visual space. These branches provide complementary language-derived priors and supervised visual evidence. The prediction module contributes to continual classification, but does not affect image--text retrieval.

We evaluate the framework under the X-TAIL protocol~\cite{xu2024rail} across ten heterogeneous image-classification datasets and assess both class-incremental learning and image--text retrieval. The experiments examine whether LoRA-NF improves over rank- and budget-matched LoRA, projected-gradient, and null-space-inspired factor variants; whether cross-modal distillation preserves retrieval after sequential adaptation; whether LR-RGDA and the CLIP text classifier provide complementary decision signals; and whether the resulting classification gains coexist with preserved cross-modal capability. \todo{Insert verified Transfer/Average/Last, classifier-ablation, and retrieval results.}

Our contributions are summarized as follows:
\begin{itemize}[leftmargin=*]
    \item We introduce \textbf{LoRA-NF}, which reparameterizes the LoRA update from $\mathbf{A}\mathbf{B}$ to $\mathbf{P}\mathbf{A}\mathbf{B}$ and persistently filters adapter inputs using historical activation spectra. We characterize its historical response, expressive equivalence to rank-matched LoRA under an invertible leaky filter, and spectrum-dependent optimization geometry.
    \item We build a complete continual-learning framework around LoRA-NF, using cross-modal knowledge distillation to preserve image--text relations and an LR-RGDA/text ensemble to combine supervised visual statistics with CLIP's semantic prior.
    \item We systematically evaluate continual classification, cross-modal retrieval, mechanism-level predictions, training and classifier ablations, and the resulting adaptation--preservation behavior under a unified protocol.
\end{itemize}
